\chapter{Appendectomy}

\section*{Introduction \& Clinical Indications}
An appendectomy is the surgical removal of the vermiform appendix, a small, finger-shaped lymphoid organ that projects from the posteromedial aspect of the cecum. This procedure is most commonly performed as an emergency intervention to treat acute appendicitis, a painful inflammation and infection of the appendix. Acute appendicitis represents one of the most prevalent abdominal emergencies encountered globally, affecting individuals across all age groups. Timely appendectomy is crucial to prevent serious and potentially life-threatening complications such as rupture of the appendix, which can lead to localized abscess formation, diffuse peritonitis, sepsis, and multi-organ failure. As such, it remains one of the most frequently performed emergency general surgical procedures worldwide, with both open and minimally invasive laparoscopic approaches widely utilized.

\begin{itemize}
  \item Appendix removal surgery
  \item Surgical appendicectomy
  \item Laparoscopic appendectomy
  \item Open appendectomy
  \item Appendicectomy
  \item Cecal appendage excision
\end{itemize}

The fundamental principle of an appendectomy is the complete and safe excision of the inflamed, infected, or diseased vermiform appendix to eliminate the source of pathology and prevent its progression. Acute appendicitis typically results from obstruction of the appendiceal lumen, often by a fecalith, lymphoid hyperplasia, foreign body, or rarely, a tumor. This obstruction leads to bacterial overgrowth, distension, inflammation, venous congestion, and eventual ischemia and necrosis of the appendiceal wall. Surgical removal halts this pathological process, thereby preventing perforation, localized abscess formation, or diffuse peritonitis, which are associated with significant morbidity and mortality.

The primary rationale for appendectomy is to definitively treat acute appendicitis and prevent its severe complications. The technical goals of the procedure include meticulous dissection to isolate the appendix from its mesoappendix, secure ligation or stapling of the appendiceal base at its junction with the cecum, and careful removal of the specimen to minimize spillage of infected contents into the peritoneal cavity. Both open and laparoscopic approaches adhere to these core principles, differing primarily in the method of abdominal access and visualization. The laparoscopic approach offers magnified views, reduced postoperative pain, and improved cosmetic outcomes, while the open approach provides direct tactile feedback and may be preferred in cases of severe inflammation, extensive adhesions, or diffuse peritonitis.

The earliest recorded successful appendectomy is attributed to Claudius Amyand, a French surgeon, who performed the procedure in 1735 on a boy with a perforated appendix within a scrotal hernia. However, it was Reginald Fitz who, in 1886, accurately described the clinical syndrome of acute appendicitis and, crucially, advocated for early surgical intervention to prevent perforation. This marked a pivotal moment in the understanding and management of the disease. Charles McBurney further advanced the understanding and surgical approach in 1889, describing the classic point of maximal tenderness (McBurney's point) and the muscle-splitting incision that bears his name, which became the standard for open appendectomy for decades.

The modern era saw a significant shift with Kurt Semm, a German gynecologist, performing the first successful laparoscopic appendectomy in 1980. This minimally invasive technique rapidly gained acceptance due to its advantages, including reduced postoperative pain, shorter hospital stays, and improved cosmetic outcomes. The widespread adoption of laparoscopy has revolutionized the surgical management of appendicitis, making it the preferred approach in most centers today, though open appendectomy remains a vital technique, particularly in complicated cases or when laparoscopic facilities are unavailable.

An appendectomy is primarily indicated for the following clinical scenarios:
\begin{itemize}
  \item \textbf{Acute Uncomplicated Appendicitis:} Characterized by inflammation of the appendix without evidence of perforation, gangrene, or abscess formation, often diagnosed based on clinical findings, laboratory tests (e.g., leukocytosis), and imaging (e.g., CT scan, ultrasound). The Alvarado Score or Pediatric Appendicitis Score can aid in risk stratification.
  \item \textbf{Acute Complicated Appendicitis:} Involving perforation, gangrenous changes, or the presence of a periappendiceal abscess or phlegmon, requiring urgent surgical intervention.
  \item \textbf{Suspected Appendicitis with Equivocal Imaging:} When clinical suspicion remains high despite inconclusive imaging studies, a diagnostic laparoscopy may be performed, often converting to appendectomy if an inflamed appendix is confirmed.
  \item \textbf{Appendiceal Mass or Phlegmon:} While sometimes initially managed non-operatively with intravenous antibiotics and percutaneous drainage, an interval appendectomy may be recommended 6-8 weeks later to prevent recurrence and rule out underlying malignancy.
  \item \textbf{Appendiceal Tumors:} Such as carcinoid tumors (most common, especially at the tip), mucinous neoplasms (e.g., mucoceles), or adenocarcinomas, which require surgical resection, sometimes with additional oncologic resections depending on tumor type and stage.
  \item \textbf{Chronic Appendicitis:} A rare and somewhat controversial diagnosis, considered when patients experience recurrent episodes of right lower quadrant pain attributed to the appendix, often after thorough exclusion of other gastrointestinal or gynecological causes.
  \item \textbf{Incidental Appendectomy:} Occasionally performed during other abdominal surgeries (e.g., gynecological procedures, colon resections) if the appendix appears abnormal or to prevent future appendicitis, though this practice is less common now.
\end{itemize}

For the vast majority of patients presenting with acute appendicitis, appendectomy is considered the primary and definitive emergency surgical treatment. Its role is to promptly remove the inflamed organ, thereby preventing progression to perforation, peritonitis, and sepsis, which are associated with increased morbidity and mortality. Early surgical intervention is the cornerstone of management for acute appendicitis.

In certain specific scenarios, appendectomy may function as a secondary or delayed procedure. For instance, in cases of a well-contained appendiceal phlegmon or abscess, initial management might involve a period of non-operative treatment with intravenous antibiotics and, if applicable, percutaneous drainage of the abscess. Following resolution of the acute inflammatory process, an "interval appendectomy" is often performed 6-8 weeks later. This secondary approach aims to prevent recurrence, allows for a safer operation in a less inflamed and vascularized field, and provides an opportunity for histological examination to rule out underlying malignancy. Furthermore, appendectomy can be an incidental finding and procedure during a diagnostic laparoscopy performed for undifferentiated abdominal pain, where an inflamed appendix is discovered unexpectedly.

\section*{Relevant Surgical Anatomy}
The vermiform appendix is a narrow, blind-ended tubular diverticulum arising from the posteromedial wall of the cecum, approximately 2-3 cm below the ileocecal valve. Its length is highly variable, typically ranging from 2 to 20 cm, with an average of 9 cm. The base of the appendix is constant, located at the convergence of the three taeniae coli of the large intestine, a crucial surgical landmark. However, its tip can assume various positions, including retrocecal (most common, 65\%), pelvic (30\%), subcecal, preileal, postileal, or paracecal, which influences the clinical presentation and surgical approach.

The appendix is suspended by a peritoneal fold known as the mesoappendix, which originates from the mesentery of the terminal ileum. This mesoappendix contains the appendicular artery, a terminal branch of the ileocolic artery (which itself arises from the superior mesenteric artery), providing the sole blood supply to the appendix. Venous drainage parallels the arterial supply, emptying into the ileocolic vein. Lymphatic drainage follows the appendicular artery to the ileocolic lymph nodes. The nerve supply to the appendix is derived from the sympathetic nervous system (T10-L1 spinal segments) and the parasympathetic vagus nerve. Visceral afferent fibers from T10 are responsible for the initial vague periumbilical pain characteristic of early appendicitis, which later localizes to the right lower quadrant as inflammation irritates the parietal peritoneum. Key surgical landmarks include:
\begin{itemize}
  \item \textbf{McBurney's Point:} Located one-third of the way from the anterior superior iliac spine to the umbilicus, often corresponding to the base of the appendix and the point of maximal tenderness in acute appendicitis.
  \item \textbf{Teniae Coli:} The three longitudinal muscle bands on the large intestine converge at the base of the appendix, serving as a reliable guide to its location.
  \item \textbf{Ileocecal Valve:} The appendix originates just inferior to this valve.
\end{itemize}

\section*{Preoperative Workup \& Preparation}
Preparation for an appendectomy is typically expedited due to the emergency nature of acute appendicitis, focusing on rapid diagnosis and stabilization. A thorough medical history and physical examination are performed to assess the patient's overall health and identify any comorbidities.

Key preoperative steps include:
\begin{itemize}
  \item \textbf{NPO Status:} Patients are instructed to refrain from eating or drinking (nil per os) immediately upon suspicion of appendicitis to minimize the risk of aspiration during general anesthesia.
  \item \textbf{Intravenous Access:} At least one large-bore intravenous line is established for fluid resuscitation, electrolyte correction, and medication administration.
  \item \textbf{Laboratory Tests:} Blood work includes a complete blood count (CBC) to check for leukocytosis and differential shift, basic metabolic panel (BMP) to assess renal function and electrolytes, coagulation studies (PT/INR, PTT), and type and screen for potential blood transfusion, especially in complicated cases. Urinalysis is performed to rule out urinary tract infection.
  \item \textbf{Imaging Studies:} Computed tomography (CT) scan of the abdomen and pelvis with intravenous contrast is the most common and accurate diagnostic tool. Ultrasound may be used, especially in children and pregnant women, to minimize radiation exposure.
  \item \textbf{Anesthesia Consultation:} An anesthesiologist evaluates the patient's overall health, comorbidities, and airway for general anesthesia, discussing potential risks and the anesthetic plan.
  \item \textbf{Prophylactic Antibiotics:} Broad-spectrum intravenous antibiotics (e.g., cefoxitin, cefazolin plus metronidazole) are administered within 60 minutes prior to incision to reduce the risk of surgical site infection and intra-abdominal abscess.
  \item \textbf{Informed Consent:} The surgical team explains the procedure, potential risks, benefits, and alternatives to the patient or their legal guardian, obtaining informed consent.
  \item \textbf{Medication Review:} Any medications, especially anticoagulants, antiplatelet agents, or immunosuppressants, are reviewed and managed appropriately in consultation with relevant specialists.
\end{itemize}

\textbf{Situations requiring stabilization or an alternative strategy:}
\begin{itemize}
  \item There are virtually no absolute contraindications to appendectomy when acute appendicitis is definitively diagnosed, particularly if complicated (e.g., perforation, gangrene). The risks of delaying surgery (e.g., progression to diffuse peritonitis, sepsis, death) far outweigh the surgical risks in such emergency situations.
\end{itemize}

\textbf{Relative Contraindications and Precautions:}
\begin{itemize}
  \item \textbf{Appendiceal Phlegmon or Abscess:} A stable, well-contained mass or abscess may be managed with antibiotics and drainage, while diffuse peritonitis or clinical deterioration requires urgent source control. Interval appendectomy is selective rather than automatic.
  \item \textbf{Severe Comorbidities:} Patients with severe cardiopulmonary disease, uncontrolled coagulopathy, or extreme frailty may require careful medical optimization before surgery. While these conditions increase surgical risk, they do not typically contraindicate an emergency appendectomy if the diagnosis is clear, but they necessitate a multidisciplinary approach.
  \item \textbf{Pregnancy:} Appendectomy in pregnant women requires careful consideration to minimize risks to both mother and fetus. The laparoscopic approach is generally favored in early trimesters, while an open approach may be preferred in later trimesters due to uterine size or when the appendix is high-riding. Fetal monitoring is often employed.
  \item \textbf{Diffuse Peritonitis:} While an indication for surgery, diffuse peritonitis may necessitate an open approach to ensure thorough abdominal lavage, address widespread contamination, and manage potential bowel injury.
  \item \textbf{Indeterminate Diagnosis with Low Clinical Suspicion:} If imaging is equivocal and clinical suspicion is low, a period of observation with serial examinations and repeat labs may be warranted. However, if suspicion remains or symptoms worsen, diagnostic laparoscopy is often performed to confirm or rule out appendicitis.
\item \textbf{Coagulopathy:} Significant bleeding disorders or patients on anticoagulation require careful management, including reversal of anticoagulation if safe, prior to surgery.
\end{itemize}

\section*{Surgical Technique \& Procedure}
An appendectomy is performed under general anesthesia. The patient is typically positioned supine on the operating table, with arms tucked or abducted, and often placed in a slight Trendelenburg position with left side down tilt for better visualization of the right lower quadrant.

\textbf{Laparoscopic Appendectomy:}
\begin{enumerate}
  \item \textbf{Access and Pneumoperitoneum:} A primary incision (usually 10-12 mm) is made at the umbilicus for insertion of a trocar and camera. Pneumoperitoneum is established by insufflating the abdomen with carbon dioxide gas to a pressure of 12-15 mmHg, creating a working space.
  \item \textbf{Trocar Placement:} Two additional working trocars (typically 5 mm) are placed under direct vision, commonly in the suprapubic region and the left lower quadrant, or right lower quadrant, depending on surgeon preference and patient anatomy.
  \item \textbf{Identification and Mobilization:} The appendix is identified, often by tracing the taeniae coli to its base. It is then carefully mobilized from surrounding structures and adhesions using atraumatic graspers.
  \item \textbf{Mesoappendix Division:} The mesoappendix, containing the appendicular artery, is meticulously divided using electrocautery, ultrasonic energy devices, or surgical clips. Care is taken to ensure complete hemostasis.
  \item \textbf{Appendiceal Base Ligation/Stapling:} The base of the appendix at its junction with the cecum is securely ligated with absorbable sutures (e.g., Endoloop) or transected using an endoscopic stapler.
  \item \textbf{Transection and Specimen Removal:} The appendix is transected distal to the ligation/staple line. The resected specimen is then placed into an endoscopic retrieval bag to prevent wound contamination and removed through one of the port sites.
  \item \textbf{Irrigation and Closure:} The abdominal cavity is inspected for hemostasis and contamination. If significant contamination is present, irrigation with saline may be performed. The abdomen is deflated, and port sites are closed in layers, with fascia closure for larger ports (>10 mm) and skin closure with sutures or surgical glue.
\end{enumerate}

\textbf{Open Appendectomy:}
\begin{enumerate}
  \item \textbf{Incision:} A single incision is made, typically a McBurney or Rocky-Davis incision (oblique or transverse muscle-splitting incision) in the right lower quadrant, designed to minimize muscle damage. In cases of diagnostic uncertainty or diffuse peritonitis, a lower midline incision may be used.
  \item \textbf{Abdominal Entry:} The abdominal wall layers (skin, subcutaneous tissue, external oblique, internal oblique, transversus abdominis, peritoneum) are incised or separated.
  \item \textbf{Identification and Mobilization:} The cecum is identified and delivered into the wound. The taeniae coli are traced to locate the appendix. The appendix is then mobilized from any surrounding adhesions.
  \item \textbf{Mesoappendix Ligation:} The mesoappendix is clamped, ligated with absorbable sutures, and divided, ensuring the appendicular artery is securely controlled.
  \item \textbf{Appendiceal Base Ligation/Transection:} The base of the appendix is crushed, ligated with a non-absorbable suture, and then transected. The appendiceal stump may be inverted into the cecum using a purse-string suture, though this practice is less common now.
  \item \textbf{Closure:} The abdominal cavity is inspected. If contamination is present, irrigation may be performed. The incision is closed in layers: peritoneum, muscle layers (if not muscle-splitting), fascia, subcutaneous tissue, and skin. Drains are generally not placed unless there is a significant abscess cavity or extensive contamination.
\end{enumerate}

\section*{Complications \& Risks}
While appendectomy is generally a safe procedure, like all surgeries, it carries potential risks. These can be broadly categorized as general surgical complications and procedure-specific risks.

\textbf{General Surgical Complications:}
\begin{itemize}
  \item \textbf{Bleeding:} Ranging from minor bruising at incision sites to significant intra-abdominal hemorrhage requiring blood transfusion or reoperation.
  \item \textbf{Infection:}
    \begin{itemize}
      \item \textbf{Surgical Site Infection (SSI):} Redness, pain, swelling, warmth, or pus at the incision site, more common in open appendectomy or perforated cases.
      \item \textbf{Intra-abdominal Abscess:} Collection of pus within the abdominal cavity, often in the pelvis or right paracolic gutter, potentially requiring percutaneous drainage or further surgery.
      \item \textbf{Peritonitis:} Widespread infection and inflammation of the abdominal lining, especially if the appendix has perforated or there is significant spillage.
    \end{itemize}
  \item \textbf{Anesthesia Complications:} Adverse reactions to medications, respiratory problems (e.g., pneumonia, atelectasis), cardiac events (e.g., arrhythmia, myocardial infarction), nausea, vomiting, or aspiration pneumonitis.
  \item \textbf{Deep Vein Thrombosis (DVT) and Pulmonary Embolism (PE):} Blood clot formation in the legs that can travel to the lungs, a rare but serious complication, mitigated by early ambulation and prophylactic measures.
\end{itemize}

\textbf{Procedure-Specific Risks:}
\begin{itemize}
  \item \textbf{Bowel Injury:} Accidental damage to the small intestine, large intestine, or cecum during dissection, particularly in cases with severe inflammation or adhesions, potentially leading to a leak, fistula, or peritonitis.
  \item \textbf{Appendiceal Stump Leak:} Failure of the ligated or stapled appendiceal stump to heal, leading to leakage of intestinal contents into the peritoneal cavity and subsequent peritonitis or abscess formation.
  \item \textbf{Adhesions and Bowel Obstruction:} Formation of scar tissue within the abdomen after surgery, which can lead to future episodes of small bowel obstruction, sometimes requiring reoperation.
  \item \textbf{Nerve Injury:} Damage to superficial nerves (e.g., ilioinguinal, iliohypogastric nerves) during incision or port placement, causing numbness, burning pain, or chronic neuralgia in the groin or thigh.
  \item \textbf{Incisional Hernia or Port-Site Hernia:} Weakness in the abdominal wall at the incision or port site, allowing abdominal contents to protrude, requiring subsequent repair. More common with larger port sites or inadequate fascial closure.
  \item \textbf{Missed Diagnosis/Negative Appendectomy:} Rarely, a normal appendix may be removed, and another cause for abdominal pain (e.g., mesenteric adenitis, Meckel's diverticulitis, gynecological pathology) is later identified. The rate of negative appendectomy should ideally be below 10-15\%.
  \item \textbf{Infertility (in women):} In cases of severe pelvic infection from a perforated appendix, extensive inflammation can theoretically lead to fallopian tube damage and subsequent infertility, though this is rare with modern surgical techniques and prompt intervention.
  \item \textbf{Retained Foreign Body:} Extremely rare, but surgical instruments, sponges, or other materials can be inadvertently left inside the abdomen.
\end{itemize}

\section*{Postoperative Care \& Recovery}
Immediately after an appendectomy, the patient is transferred to the Post-Anesthesia Care Unit (PACU) for close monitoring of vital signs, pain levels, and recovery from general anesthesia. Pain management is initiated, typically with intravenous analgesics, transitioning to oral medications as tolerated. Nausea and vomiting are common in the immediate postoperative period and are managed with antiemetics. Early ambulation is strongly encouraged to prevent complications such as deep vein thrombosis, pulmonary embolism, and to promote the return of normal bowel function.

Within the first 24-48 hours, patients are gradually advanced from clear liquids to a regular diet as bowel sounds return, flatus is passed, and nausea subsides. Most patients undergoing uncomplicated laparoscopic appendectomy are discharged within 24-48 hours, provided their pain is well-controlled with oral medication, they are tolerating oral intake, and they are ambulating independently. Patients with complicated appendicitis (e.g., perforation, abscess) may require a longer hospital stay for continued intravenous antibiotics, more intensive pain management, and observation for potential complications. Full recovery typically takes 1-2 weeks for laparoscopic cases and 2-4 weeks for open cases, with patients advised to avoid heavy lifting and strenuous activities for several weeks to allow for proper healing of the abdominal wall and minimize the risk of hernia formation.

During an appendectomy, the surgeon meticulously documents the intraoperative findings, which include the macroscopic appearance of the appendix (e.g., inflamed, edematous, gangrenous, perforated), the presence of periappendiceal fluid or pus, the condition of surrounding organs, and any adhesions. This documentation is crucial for understanding the severity of the disease and guiding postoperative management. The removed appendix is routinely sent to a pathology laboratory for microscopic examination.

The pathology report is crucial for confirming the diagnosis of acute appendicitis and identifying its specific type (e.g., catarrhal, phlegmonous, gangrenous, perforated). It details the presence and extent of inflammation, the presence of fecaliths, and any signs of necrosis or perforation. Importantly, the pathologist will also rule out other conditions that can mimic appendicitis, such as appendiceal tumors (e.g., carcinoid tumors, mucinous neoplasms, or adenocarcinomas), which, though rare, have significant implications for prognosis and further treatment. The final pathology report provides definitive histological confirmation of the diagnosis and ensures no unexpected pathologies were missed.

A normal and successful postoperative course following an appendectomy is characterized by a rapid and progressive resolution of the patient's preoperative symptoms. Abdominal pain, which was the primary complaint, should significantly diminish within the first 24-48 hours, transitioning from severe to mild discomfort manageable with oral analgesics. Any preoperative fever or leukocytosis should resolve promptly. Patients typically experience a gradual return of normal bowel function, with the passage of flatus and bowel movements within 1-3 days. Wound healing should proceed without signs of infection, such as increasing redness, swelling, warmth, or purulent discharge. Patients should be able to tolerate a regular diet within a few days and progressively increase their activity levels, returning to most normal daily activities within 1-2 weeks for laparoscopic surgery and 2-4 weeks for open surgery, without significant pain or limitations.

Post-appendectomy care for uncomplicated cases is generally straightforward, focusing on wound healing, pain management, and gradual return to normal activities.
\begin{itemize}
  \item \textbf{Post-operative Visit:} A follow-up appointment is typically scheduled 1-2 weeks after surgery for a wound check, assessment of overall recovery, and to address any patient concerns.
  \item \textbf{Suture/Staple Removal:} If non-absorbable sutures or staples were used for skin closure, they will be removed at the follow-up visit. Absorbable sutures do not require removal.
  \item \textbf{Activity Restrictions:} Patients are advised to gradually increase activity. They should avoid heavy lifting (typically >10-15 lbs) and strenuous abdominal exercise for 2-4 weeks after laparoscopic surgery and 4-6 weeks after open surgery to prevent incisional hernia formation. Light activities and walking are encouraged.
  \item \textbf{Diet:} No specific dietary restrictions are usually needed once a regular diet is tolerated, though some patients may prefer bland foods initially.
  \item \textbf{Wound Care:} Patients are instructed to keep the incision sites clean and dry. Showers are generally permitted after 24-48 hours, but baths or swimming should be avoided until wounds are fully healed.
  \item \textbf{Pain Management:} Oral analgesics are prescribed and should be taken as needed, with instructions on tapering their use as pain subsides.
\item \textbf{Warning Signs:} Patients are instructed to contact their surgeon immediately if they experience persistent or worsening abdominal pain, fever (temperature >101.5$^{\circ}$F or 38.6$^{\circ}$C), chills, persistent nausea or vomiting, inability to tolerate oral intake, increasing redness, swelling, warmth, or pus drainage from the incision site, or any signs of bowel obstruction.
\end{itemize}
Routine imaging or laboratory tests are generally not required unless complications are suspected based on clinical presentation.

\section*{Outcomes \& Prognosis}
Acute appendicitis is one of the most common causes of acute abdominal pain requiring emergency surgery, with a lifetime risk estimated to be between 7\% and 10\%. Its incidence peaks in the second and third decades of life (ages 10-30), although it can occur at any age, from infancy to old age, with diagnostic challenges often greater in the very young and elderly. There is a slight male predominance, with a male-to-female ratio of approximately 1.4:1. The most common indication for appendectomy is acute uncomplicated appendicitis, accounting for the vast majority of cases. Mortality rates for uncomplicated appendicitis are very low, typically less than 0.1\%. However, mortality can rise significantly in cases of perforated appendicitis, especially in the elderly, immunocompromised individuals, or those with significant comorbidities, reaching up to 5\% or higher. Overall morbidity, including complications like wound infection and intra-abdominal abscess, ranges from 5\% to 10\%, with higher rates observed in complicated cases and in patients undergoing open appendectomy.

The outcomes following an appendectomy for acute appendicitis are generally excellent, with high success rates and a favorable prognosis, particularly when the diagnosis is made early and the appendix has not perforated. For uncomplicated appendicitis, the success rate of appendectomy in resolving symptoms and preventing complications exceeds 95\%. Patients typically experience a full functional recovery, returning to their normal activities within a few weeks. Survival outcomes are also very high, especially in younger, otherwise healthy individuals. Recurrence is not an issue as the appendix is completely removed. Prognostic factors influencing outcomes include the stage of appendicitis at presentation (uncomplicated vs. complicated with perforation or abscess), patient age, the presence of comorbidities, and the surgical approach. Patients with perforated appendicitis or an appendiceal abscess have a longer hospital stay, higher rates of postoperative complications (e.g., intra-abdominal abscess, wound infection), and a more prolonged recovery period, but the long-term prognosis remains good for most. Landmark trials like the APPAC trial have explored non-operative management with antibiotics for uncomplicated appendicitis, demonstrating its feasibility in selected cases, though appendectomy remains the definitive treatment and gold standard.

\section*{References}
\begin{enumerate}
  \item MedlinePlus. Appendectomy. U.S. National Library of Medicine; 2024. Available from: \url{https://medlineplus.gov/appendectomy.html}
  \item Brunicardi FC, Andersen DK, Billiar TR, et al., editors. Schwartz's Principles of Surgery. 11th ed. McGraw-Hill Education; 2019.
  \item Society of American Gastrointestinal and Endoscopic Surgeons (SAGES). Guidelines for Laparoscopic Appendectomy. Available from: \url{https://www.sages.org/publications/guidelines/guidelines-laparoscopic-appendectomy/}
  \item Sabiston DC, Townsend CM, Beauchamp RD, et al., editors. Sabiston Textbook of Surgery: The Biological Basis of Modern Surgical Practice. 21st ed. Elsevier; 2021.
\end{enumerate}

\clearpage
